\section{Arithmetic for the signed determinant}
\label{sec:v-arithmetic}

This section supplies the arithmetic interface for the negative
\(V\)-family.  The base is fixed before \(m\) tends to infinity.  Thus the
statement is made separately for \(q=2,3,4,5,6,7\), with \(a=-q\); it is
not a theorem uniform in an unbounded parameter \(q\).  The proof has four
parts: an exact polynomial divisor, residue-block valuations away from the
fixed primes, an outer pole assignment, and a weighted prime-number-theorem
sum.  All subsequent limits follow from the displayed residue valuations and
outer assignment bounds.

Put \(m=4n\), \(s=k=h=n\), and
\[
 I=[2n,6n),\qquad J=[n,10n),\qquad
 \Lambda=\{1,\ldots,n\},\qquad N=10n-1,\quad U=11n-1.
\]
For an integer parameter \(a\), define
\[
 t_j=\sum_{r=1}^{j}\frac{a^{j-r}}r,\qquad
 b_j=\sum_{r=1}^{j}\frac{a^{j-r}}{r^2},
 \qquad t_0=b_0=0,
\]
and, for \(i\in I\), \(j\in J\), \(\ell\in\Lambda\),
\[
 C_{ij}=\begin{cases}(a^i-a^j)/(j-i),&i\ne j,\\0,&i=j,\end{cases}
 \quad
 B_{ij}=\begin{cases}(t_i-t_j)/(j-i),&i\ne j,\\b_i,&i=j,\end{cases}
\]
\[
 A_{ij}=a^i\delta_{ij},\qquad V_{\ell j}=\frac1{j+\ell},\qquad
 G_m(T;a)=\det\begin{bmatrix}C(a)\\ B(a)-TA(a)\\ V\end{bmatrix}.
\]
The determinant has size \(2m+n\).  All valuations below are coefficientwise
in \(T\): \(v_p(F)=\min_r v_p([T^r]F)\).

\subsection{Denominator and the formal divisor}

Let \(d_r=\operatorname{lcm}(1,\ldots,r)\) and
\[
 \Delta_m=d_N^{3m}d_U^n,\qquad
 D_m=\left\lceil\frac{53m^2}{48}\right\rceil-m
      =\left\lceil\frac{53n^2}{3}\right\rceil-4n.
\]

\begin{lemma}[denominator]
\label{lem:v-denom}
For an indeterminate \(a\),
\[
 \Delta_mG_m(T;a)\in\mathbb Z[a,T].
 \]
\end{lemma}

\begin{proof}
Every denominator in a \(C\)-entry divides \(d_N\).  In an off-diagonal
\(B\)-entry one harmonic denominator and one divided-difference denominator
occur, while \(b_i\) has two harmonic denominators; hence \(d_N^2\) clears
each lower row.  The denominators of the \(V\)-rows divide \(d_U\).
Multiplying the rows and expanding the determinant proves the assertion.
\end{proof}

\begin{lemma}[formal $a$-divisor]
\label{lem:v-formal-divisor}
Coefficientwise in \(T\),
\[
 \Delta_mG_m(T;a)\in a^{D_m}(a-1)^m\mathbb Z[a,T].
\]
\end{lemma}

\begin{proof}
The following argument is coefficientwise in \(\mathbb Q[[a]][T]\), so
every use of an infinite series can be replaced by a finite truncation for a
fixed coefficient.  For \(u_t(r)=(r-t)^{-1}\) when \(r\ne t\), and
\(u_t(t)=0\), put \(S_1(a)=\sum_{v\ge1}a^v/v\) and
\(S_2(a)=\sum_{v\ge1}a^v/v^2\).  Direct coefficient comparison, including
the diagonal entries, gives
\[
 C_{ij}=a^i u_i(j)+a^j u_j(i),
\]
\[
 (B-S_1C)_{ij}
 =\sum_{t\ge0}a^tu_t(i)u_t(j)-\frac{a^i+a^j}{(i-j)^2}
 \quad(i\ne j),
\]
and
\[
 (B-TA)_{ii}=\sum_{t\ge0}a^tu_t(i)^2-a^i(T+S_2(a)).
\]
The last formula is the diagonal term and is the reason no off-diagonal
limiting convention is needed.

Expand first along the \(n\) rows of \(V\).  In each remaining \(2m\)-row
minor, use the preceding rank-one expansion.  A nonzero term selects
\(2m=8n\) distinct features.  If \(b,x,y,z\) count, respectively, the
series features, the upper boundary features, the lower boundary features,
and the column boundary features, then \(b+x+y+z=8n\).  Distinctness gives
the lower bound
\[
 \Psi\ge {b\choose2}+{x\choose2}+x\max(2n,b)
       +2ny+{y\choose2}+nz+{z\choose2}.
\]
Here the \(b\) series indices are distinct; an upper boundary index cannot
be one of them because the two corresponding right vectors coincide.  The
\(x\) and \(y\) indices are distinct members of \(I=[2n,6n)\), while the
\(z\) column indices are distinct members of \(J=[n,10n)\).  These are the
only features used in the count, so a collision can only delete a term and
cannot lower the displayed lower bound.
Writing \(B_0=b+x\), the right side is at least
\[
 Q(B_0,y)=B_0^2+B_0y+y^2-9nB_0-7ny+40n^2-4n
 =\frac{53n^2}{3}-4n+r^2+rs+s^2
\]
for suitable real \(r,s\).  Since the order is integral, every coefficient
of every \(2m\)-minor has order at least \(D_m\).  The \(V\)-minors are
rational constants independent of \(a,T\), so expansion along \(V\) gives
\(a^{D_m}\mid G_m(T;a)\) in \(\mathbb Q[a,T]\).  Combining this with the
denominator lemma and monic division gives
\(a^{D_m}\mid\Delta_mG_m\) in \(\mathbb Z[a,T]\).

Finally every \(C\)-entry is divisible by \(a-1\) in \(\mathbb Q[a]\),
including the zero diagonal convention.  There are \(m\) such rows, hence
\((a-1)^m\mid G_m\).  The two monic factors \(a^{D_m}\) and \((a-1)^m\)
are coprime, so their product divides the integral polynomial
\(\Delta_mG_m\).  This proves the lemma for every \(n\ge1\), not only for
the finite verification values.
\end{proof}

For \(a=-q\), the formal divisor implies the exact coefficientwise bound
\[
 v_p([T^r]G_m(T;-q))\ge v_p(q)D_m+v_p(q+1)m
       -3m\lfloor\log_pN\rfloor-n\lfloor\log_pU\rfloor .
\]
The two credited prime sets are disjoint because \(\gcd(q,q+1)=1\).  Notice
that this argument includes composite \(q\) and powers \(p^e\); no
square-free assumption is being made.
For reference, the finite exceptional sets in the five middle cases are
\[
\begin{array}{c|c@{\qquad}c|c}
q&\mathcal P_q=\{p:p\mid q(q+1)\}&q&\mathcal P_q\\ \hline
3&\{2,3\}&4&\{2,5\}\\
5&\{2,3,5\}&6&\{2,3,7\}\\
7&\{2,7\}&&
\end{array}
\]
The \(q=2\) row is \(\{2,3\}\), and is treated by the same rule.

\subsection{The $p$-unit residue blocks}

Fix an odd prime \(p\nmid q(q+1)\) with \(p^2>\max(N,U)\), and write
\(i=u+kp\), \(j=u+\ell p\), \(0\le u<p\).  Since \(a=-q\) and
\(a^p\equiv a\pmod p\), splitting the defining sums at multiples of \(p\)
gives
\begin{align*}
 p t_{u+kp}
 &=\sum_{\substack{1\le r\le u+kp\\p\nmid r}}
       \frac{p a^{u+kp-r}}r
   +\sum_{s=1}^{k}\frac{a^{u+p(k-s)}}s
   \equiv a^u t_k\pmod p,\\
 p^2 b_{u+kp}
 &=\sum_{\substack{1\le r\le u+kp\\p\nmid r}}
       \frac{p^2a^{u+kp-r}}{r^2}
   +\sum_{s=1}^{k}\frac{a^{u+p(k-s)}}{s^2}
   \equiv a^u b_k\pmod p,\qquad
 p^2TA_{ij}\equiv0\pmod p.
\end{align*}
Here \(k,\ell<p\) follows from \(i\le N\), \(j\le U\), and
\(p^2>\max(N,U)\), so each displayed denominator \(s\) is a \(p\)-unit.
Thus the scaled entries satisfy
\[
 pC_{u+kp,u+\ell p}\equiv a^uC_{k\ell}\pmod p,\qquad
 p^2B_{u+kp,u+\ell p}\equiv a^uB_{k\ell}\pmod p,
\]
whereas
\[
 pV_{r,u+\ell p}\equiv0\quad\text{unless }r+u=pw,\qquad
 pV_{r,u+\ell p}=\frac1{\ell+w}\quad(r+u=pw).
\]
The denominators in the last expression are units because \(p^2>U\).
Cross-residue \(C\) and \(B\) entries vanish after the displayed row
scalings.  The reduced residue-\(u\) block has \(2|K_u|+|W_u|\) rows and
\(|L_u|\) columns.  Therefore its rank over \(\mathbb F_p\) is at most
\[
 \min(2|K_u|+|W_u|,|L_u|)
\]
where
\[
 K_u=\{k:1/2\le u/m+kp/m<3/2\},\quad
 L_u=\{\ell:1/4\le u/m+\ell p/m<5/2\},
\]
\[
W_u=\{w\ge1:0<w p/m-u/m\le1/4\}.
\]
Let $\operatorname{rk}_u$ denote the actual rank of this reduced block.
After permuting rows and columns by residue class, the globally scaled square
matrix of order $2m+n$ is block diagonal with precisely these blocks.  Hence
its rank modulo $p$ equals $\sum_u\operatorname{rk}_u$, and its determinant
has valuation at least
\[
 (2m+n)-\sum_u\operatorname{rk}_u.
\]
The total row-scaling factor is $p^{3m+n}$: the $C$, $B$, and $V$ rows were
scaled by $p$, $p^2$, and $p$, respectively.  Undoing this factor gives
\[
 v_p(G_m(T;-q))\ge -\left(m+\sum_u\operatorname{rk}_u\right).
\]
These valuations are coefficientwise in $T$.  Indeed, the Gauss valuation on
$\mathbb Q_p(T)$ is discrete with residue field $\mathbb F_p(T)$; Smith
normal form over its valuation ring gives
$v_p(\det M)\ge (2m+n)-\operatorname{rank}(M\bmod p)$ for every integral
square matrix $M$ of order $2m+n$.  Thus only the actual block ranks
$\operatorname{rk}_u$ are credited,
not a count of non-divisible rows.  Summing the block ranks gives the
exponent used below.
The only rank saving deliberately credited in the valuation sum is the profile
\(K=\{0\},L=\{0,1\},W=\varnothing\).  After dividing its rows by the
unit \(a^u\), its quotient matrix is
\[
 \begin{pmatrix}0&1-a\\0&-1\end{pmatrix},
\]
which has rank one over every residue field.  This is a universal identity;
no accidental rank deficiency of another profile is credited.

For \(\rho=p/m\), define
\[
 K(x,\rho)=\{k\ge0:1/2\le x+k\rho<3/2\},\quad
 L(x,\rho)=\{\ell\ge0:1/4\le x+\ell\rho<5/2\},
\]
\[
 W(x,\rho)=\{w\ge1:0<w\rho-x\le1/4\},
\qquad 0\le x<\rho,
\]
and define
\[
\begin{aligned}
 \mathfrak r(x,\rho)
   &=\min(2|K(x,\rho)|+|W(x,\rho)|,|L(x,\rho)|)\\
   &\quad-\mathbf 1_{\{K(x,\rho)=\{0\},\,
       L(x,\rho)=\{0,1\},\,W(x,\rho)=\varnothing\}},\\
 \mathfrak c(\rho)&=1+\int_0^\rho\mathfrak r(x,\rho)\,dx .
\end{aligned}
\]
with the rank-one profile subtracted on its exact active cells.  The endpoint
conventions are left-closed/right-open for \(K,L\), and
\(w\rho-1/4\le x<w\rho\) for \(W\).  Every wall has the form
\(x=c+b\rho\) with \(4c,b\in\mathbb Z\).  Since \(m=4n\), the function is
constant on each lattice interval \([u/m,(u+1)/m)\), and hence
\[
 \sum_{u=0}^{p-1}\mathfrak r(u/m,p/m)
   =m\int_0^{p/m}\mathfrak r(x,p/m)\,dx.
\]
The residue-block lemma is therefore
\[
 v_p(G_m(T;-q))\ge -m\mathfrak c(p/m).
\]
The proof uses only the p-unit conditions and the index geometry, so it is
independent of the value of the fixed \(q\).

For the arrangement with cutoff $R=40$, all active walls on
\(1/40<\rho<5/4\) are enumerated by the finite affine rule above.  Define
the finite critical set explicitly by
\[
 \mathcal W_{40}=\{1/40,5/4\}\cup
 \left\{\frac{c-c'}{b'-b}:\ x=c+b\rho,\ x=c'+b'\rho
 \text{ are listed walls, }\frac1{40}<\frac{c-c'}{b'-b}<\frac54\right\}.
\]
At each open interval of this set, choose the least rational midpoint and
generate the $x$-cells and profiles by the endpoint inequalities above;
the profile is constant on the resulting cells.  Their
exact ordering gives 412 \(\rho\)-slabs, 2472 \(x\)-cells and 582 profiles;
the one intentionally deficient profile integrates to the saving \(5/48\).
More explicitly, retain on \(1/40\le\rho\le5/4\) and \(0\le x\le\rho\)
the affine walls
\[
\begin{gathered}
x=0,\quad x=\rho,\quad
x=\tfrac12-k\rho,\quad x=\tfrac32-k\rho,\\
x=\tfrac14-\ell\rho,\quad x=\tfrac52-\ell\rho,\quad
x=w\rho-\tfrac14,\quad x=w\rho .
\end{gathered}
\]
for all nonnegative \(k,\ell\) and positive \(w\) whose line meets the
closed strip.  Only finitely many such lines occur; sort their pairwise
rational intersections and evaluate the exact midpoint profile on each cell.
This finite rule gives the 412 \(\rho\)-slabs, 2472 \(x\)-cells and 582
profiles without floating-point geometry.
The corresponding exact finite comparisons use only the displayed residue
blocks and the stated rank formula.

\subsection{Outer poles and precedence}

For \(p>5n\), use row potentials \(\lambda_B(i),\lambda_V(\ell)\) and
column potentials \(\chi(j)\) as follows:
\[
\begin{array}{c|c|c|c}
 \text{range}&(\lambda_B,\lambda_V,\chi)&\text{cost}&\text{pole}\\ \hline
 5n<p<6n&(1+[i\ge p],1,[j\ge p+2n])&19n-2p&C,B,V\\
 6n\le p<8n&([i+p<10n],1,[j\ge p])&19n-2p&C,B,V\\
 8n\le p<10n&(0,1,[j\ge p])&11n-p&C,B,V\\
 10n\le p\le U&(0,[p-\ell\in J],0)&11n-p&V
\end{array}
\]
Here brackets are indicator functions.  A \(C\)-entry has at most a simple
pole at \(j=i+p\); a \(B\)-entry has the simple harmonic pole and, when the
residue agrees, the second pole (including the diagonal double pole); and a
\(V\)-entry has a simple pole at \(j+\ell=p\) or \(2p\).  The potentials in
the table dominate these orders entry by entry.  In particular, the two
families of \(V\)-poles are explicitly included, and no pole is assigned to
two ranges.  This proves
\[
 v_p(G_m(T;-q))\ge-E_p^{\rm out},\qquad
 E_p^{\rm out}=\begin{cases}19n-2p,&5n<p<8n,\\11n-p,&8n\le p\le U.
\end{cases}
\]
The \(C\)-row potential in the last row is \(0\); its absence there records
that no \(C\)-pole occurs once \(p>10n\).
The argument uses only p-unit cancellation and pole locations; it is unchanged
when \(-2\) is replaced by any fixed \(-q\) in the p-unit branch.

For completeness, fix the finite set \(\mathcal P_q=\{p:p\mid q(q+1)\}\).  Let
\(\mathcal W_{40}\) be the exact set of affine \(\rho\)-walls in this
arrangement, with \(1/40\) and \(5/4\) included as closed endpoints.  Define
the baseline exponent
\[
 e_{\rm nat}(p)=3m\left\lfloor\log_p N\right\rfloor
                 +n\left\lfloor\log_p U\right\rfloor,
\]
which is the coefficientwise exponent supplied by \(\Delta_m=d_N^{3m}d_U^n\),
and assign the exponent to each prime by the following ordered rule:
\[
 E_{p,q}=\begin{cases}
 e_{\rm nat}(p)-v_p(q)D_m-v_p(q+1)m,&p\in\mathcal P_q,\\[2pt]
 e_{\rm nat}(p),&p\notin\mathcal P_q,\ p\le m/40,\\
 e_{\rm nat}(p),&p\notin\mathcal P_q,\ p^2\le\max(N,U),\\
 e_{\rm nat}(p),&p\notin\mathcal P_q,\ p/m\in\mathcal W_{40},\\[2pt]
 m+\sum_u\mathfrak r(u/m,p/m),&
 \substack{p\notin\mathcal P_q,\ 1/40<p/m<5/4,\\ p/m\notin\mathcal W_{40}},\\[2pt]
 E_p^{\rm out},&5n<p\le U,\ p\notin\mathcal P_q,\\[2pt]
 0,&p>U.
 \end{cases}
\]
The first matching line always takes precedence; the cases are therefore an
ordered partition even when a small-prime condition overlaps a wall.
The outer line begins after the inner range.
The coefficientwise bounds above imply
\[
 D^*_{m,q}:=\prod_p p^{E_{p,q}},\qquad
 D^*_{m,q}G_m(T;-q)\in\mathbb Z[T].
\]
The exponent in the first line can be negative; this is allowed because the
formal divisor supplies the corresponding coefficient valuation.  No primitive
content or gcd equality is asserted.

\subsection{Prime summation and the fixed-base rate}

Let \(\vartheta(y)=\sum_{p\le y}\log p\) and
\[
 \eps_{\rm PNT}(X)=\sup_{y\ge X}\left|\frac{\vartheta(y)}y-1\right|.
\]
For every affine profile \(f\) on a rational \(\rho\)-cell,
Stieltjes summation gives
\[
 m\sum_{t_0m<p<t_1m}f(p/m)\log p
 =m^2\int_{t_0}^{t_1}f(\rho)\,d\rho
 +O\!\left(\eps_{\rm PNT}(m/40)m^2+m\log(3m)\right).
\]
The baseline strip contributes \(13/160\), the deliberately credited profile
contributes \(-5/48\), and the outer potentials contribute \(45/32\).  The
exact rational value for this arrangement (including all intervening cells) is
\[
 \begin{aligned}
 A_{40}&=N_{40}/D_{40}=5.082993964478118\ldots,\\
 N_{40}&=132263363803421925450338622847809742317124551497,\\
 D_{40}&=26020759561732371877438650035601516684254668800.
 \end{aligned}
\]
Higher prime powers contribute \(O(m^{3/2}\log(3m))\).  The exact sum over
the fixed exceptional set is
\[
 \sum_{p\in\mathcal P_q}E_{p,q}\log p
 =\sum_{p\in\mathcal P_q}e_{\rm nat}(p)\log p-D_m\log q-m\log(q+1).
\]
Only the first term on the right is the exceptional-prime cost; since
\(\mathcal P_q\) is fixed and \(N,U=O(m)\), it is \(O_q(m\log(3m))\).  The
\((q+1)^m\) formal factor contributes the separate linear term
\(-m\log(q+1)\).  As
\(D_m=53m^2/48-m+O(1)\), the other formal term contributes
\(-\frac{53}{48}m^2\log q+O_q(m)\).  Thus
\[
 \log D^*_{m,q}\le
 \left(A_{40}-\frac{53}{48}\log q\right)m^2+r_{q,m},
\]
where
\[
 r_{q,m}=O\!\left(\eps_{\rm PNT}(m/40)m^2
 +m^{3/2}\log(3m)+m\log(3m)+C_q^{\rm exc}m\log(3m)+B_qm\right),
\]
\(C_q^{\rm exc}=13|\mathcal P_q|/4\), and \(B_q\) is fixed after \(q\) is fixed.  Since
\(\eps_{\rm PNT}(X)\to0\), this is \(o_q(m^2)\).  The quadratic rates are
\[
\begin{array}{c|c}
q&A_{40}-(53/48)\log q\\\hline
3&3.869942895741\\
4&3.552293940742\\
5&3.305906269499\\
6&3.104592883873\\
7&2.934384841563
\end{array}
\]
For \(q=2\), the same formal rate is
\(4.317643952609845\ldots\); its analytic passage is the \(q=2\) instance of
Section~\ref{sec:v-analytic}.

\begin{proposition}[fixed-base arithmetic transfer]
\label{prop:v-fixed-arithmetic}
For each fixed \(q\in\{2,3,4,5,6,7\}\), the determinant has a rational
coefficientwise integerizer satisfying
\[
 \limsup_{m\to\infty,\ 4\mid m}m^{-2}\log D^*_{m,q}
 \le A_{40}-\frac{53}{48}\log q.
\]
For \(q=2,3,4,5,6,7\), the proof is the formal divisor, residue/outer bounds,
wall precedence and PNT argument above.  The finite comparisons instantiate
the displayed formulas; they do not replace the all-$m$ argument.
\end{proposition}
